\begin{abstract}
This article examines companion processes that reproduce the exact global growth of 2-gaps but change where each filter removes them. Every parent produces $r$ copies and exactly two are removed, as in the sieve sequence. The companion may place those two removals randomly, protect a chosen target, or direct them toward it. This construction separates the number of surviving 2-gaps from their location near the head.
A filter's local destruction fraction $f_r$ is compared with the random rate $2/r$ through $w_r=rf_r/2$. Under the stated blind-placement premise, the cumulative product proves that every fixed finite value of $w_r$ preserves square-window 2-gaps; with the additional availability and mixing conditions, it produces head 2-gaps infinitely often. The first boundary occurs when $w_r=1+c\log r$: square windows survive for $c < 1$, while head recurrence survives for $c < 1/2$. Under the same spatial premises, a separate exact-quota random-location model gives the same local frontiers when its normalized quotas satisfy the CRT-rate cumulative sum and summable finite-population error conditions derived below. It preserves an accepted-strike count, rather than the balanced per-parent recurrence; retaining one CRT strike count alone is insufficient. None of the spatial, availability, or mixing premises used by these conditional theorems is proved for the real sieve. Therefore, proving that the real sieve remains below the head frontier, together with persistent availability and the deterministic discrepancy bound stated in Section~10, would establish the twin-prime conjecture.
\end{abstract}
